% Generated draft: requires mathematical and source review.
\documentclass[11pt]{article}
\usepackage[margin=25mm]{geometry}
\usepackage{fontspec}
\setmainfont{Latin Modern Roman}
\usepackage{amsmath,amssymb,mathtools,unicode-math}
\setmathfont{Latin Modern Math}
\usepackage{enumitem,needspace,xurl,hyperref}
\setlength{\parindent}{0pt}
\setlength{\parskip}{5pt}
\setlength{\emergencystretch}{3em}
\newcommand{\N}{\mathbb{N}}
\newcommand{\Z}{\mathbb{Z}}
\newcommand{\Q}{\mathbb{Q}}
\newcommand{\R}{\mathbb{R}}
\newcommand{\C}{\mathbb{C}}
\newcommand{\F}{\mathbb{F}}
\newcommand{\A}{\mathbb{A}}
\newcommand{\PP}{\mathbb{P}}
\newcommand{\E}{\mathbb{E}}
\newcommand{\eps}{\varepsilon}
\newcommand{\dd}{\,\mathrm{d}}
\begin{document}
\section*{UnsolvedMath ID: 8800040}
\textbf{Source problem code:} catalogue entry
\subsection{Definitions and mathematical statement}
Let $E/\Q$ be an elliptic curve and interpret torsion free as $E(\Q)_{\mathrm{tors}}=\{O\}$. This does not mean $E(\Q)$ has rank zero. At a prime $p$ of good reduction, let $N_p=\#E(\F_p)$ and define
\[
\pi_E(x)=\#\{p\leq x:p\text{ prime of good reduction and }N_p\text{ prime}\}.
\]
The literal source asks whether every such torsion-free $E$ satisfies $\pi_E(x)\to\infty$ as $x\to\infty$.

The source summary relates this to Koblitz's conjecture and its refined local/entanglement obstructions. Trivial rational torsion alone does not state that all divisibility obstructions in the isogeny/Galois representation data have been excluded. The summary reports curves with no prime orders, but does not identify a curve or verify which of the source's torsion/admissibility hypotheses it meets. Therefore the universal literal assertion and the intended unobstructed version must be checked separately.

In a refined asymptotic formulation one expects, for suitable curves with positive conjectural constant,
\[
\pi_E(x)\sim C_E\frac{x}{(\log x)^2},\qquad C_E>0.
\]
The record supplies no definition of $C_E$ or exact admissibility criterion. Positivity cannot simply be assumed from the word torsion-free, and the original request is infinitude, weaker than this asymptotic. Results averaged over curves and counts of almost-prime group orders do not establish prime-order infinitude for a specified fixed curve.
\subsection{Short English statement}
For a fixed rational elliptic curve without rational torsion, are there infinitely many good primes with prime group order? The refined conjecture needs a precise absence-of-obstructions condition, which the source does not provide.
\subsection{Sources}
\begin{itemize}
\item Local UnsolvedMath registry, permanent ID 8800040, exact source code catalogue entry, ``Prime reductions of elliptic curves over the rationals''; attributed to Igor Shparlinski, no proposed year supplied. Registry status open, classification OPEN-TRIAGE.
\item \emph{IPAM Workshop on Number Theory and Cryptography---Open Problems} (2006), attachment 6224, PDF page 20, ``Holy Grail'', \url{https://www.ipam.ucla.edu/programs/workshops/workshop-i-number-theory-and-cryptography-open-problems/?tab=schedule}. Supplied visual audit accessed July 29, 2026; provenance records 17 PDFs/two legacy PPT decks inspected with repeated builds excluded. No PDF/PPT was accessed here.
\item The summary names Koblitz (1988), Jones for an obstruction/counterexample to the original constant, and Zywina, ``A refinement of Koblitz's conjecture'', for a corrected constant and number-field extension, without exact references or the relevant curve. No automatic claim that every source-torsion-free curve is admissible is inferred.
\item The summary names Balog--Cojocaru--David for average results and Miri--Murty, Steuding--Weng, David--Wu and Cojocaru for sieve/almost-prime progress, without precise hypotheses or citations. Its wording about almost-prime $p$ is understood only as an informal progress description; prime-order counting and almost-prime orders remain distinct.
\item Provenance flags current primary-literature status and redistribution terms NEEDS\_REVIEW; the supplied assessment says fixed-curve prime infinitude is unresolved in the intended setting.
\end{itemize}
\end{document}
